\documentclass[10pt]{article}
\title{The Promise and Peril of Cross Grade Ability Grouping \\ Evidence from Kenya and Liberia}
\author{Guthrie Gray-Lobe}
\date{}
%\usepackage[utf8]{inputenc}
%\usepackage[margin=1 in]{geometry}
\usepackage[letterpaper,top=1in,right=1in,left=1in,bottom=1in]{geometry}
\usepackage{amsmath}
\usepackage{amsfonts}
\usepackage{amssymb}
\usepackage{graphics}
\usepackage{graphicx}
\usepackage{placeins}
\usepackage{float}
\usepackage{subfigure}
\usepackage{lscape}
\usepackage{pdflscape}
\usepackage{setspace}
\usepackage{standalone}
\usepackage{fancyhdr}
\usepackage{framed}
%\pagestyle{fancy}
\usepackage{color}
\usepackage{tikz}
\usepackage{pgfplots}
\usepackage{longtable}
\usepackage{booktabs,caption}
\usepackage[draft]{todonotes}
\usepackage{booktabs,caption}
\usepackage[flushleft]{threeparttable}
\usepackage{bbm}
\usepackage{multirow}
\usepackage{lineno}
\usepackage{titlesec}
\usepackage{hyperref}
\usepackage{natbib}
\bibliographystyle{unsrtnat}

\setcounter{secnumdepth}{4}
\titleformat{\paragraph}
{\normalfont\normalsize\bfseries}{\theparagraph}{1em}{}
\titlespacing*{\paragraph}
{0pt}{3.25ex plus 1ex minus .2ex}{1.5ex plus .2ex}

\pgfplotsset{compat=1.15}
\newtheorem{theorem}{Theorem}[section]
\newtheorem{assumption}{Assumption}[section]


\linenumbers 

\begin{document}


\begin{titlepage}
\author{Guthrie Gray-Lobe\footnote{University of Chicago, email: \protect\href{mailto:graylobe@uchicago.edu }{graylobe@uchicago.edu}} \and Mridul Joshi \and Michael Kremer\footnote{University of Chicago, email: \protect\href{mailto:kremer.m@gmail.com}{kremer.m@gmail.com}} \and Joost de Laat \footnote{Utrecht University, email: \protect\href{mailto: joostdelaat@gmail.com}{joostdelaat@gmail.com}}}

\date{PLEASE DO NOT CIRCULATE \\ \today \\}
\maketitle

\begin{abstract}
This document just lays out a simple empirical framework for thinking of the cross-grade ability grouping A/B tests.
\end{abstract}
\setcounter{page}{0}
\thispagestyle{empty}
\end{titlepage}
\doublespacing

\section{Empirical Framework}

Consider a problem where we are trying to classify kids into two groups. Students have characteristics $X_i$ and a running score $r_i$ that will be used to classify them. Rather than think of this in terms of students being ``switched'' between groups, it's probably going to be easier if we are just choosing whether to place some students in a ``high'' level course from a default low level course. 

The running score $r_i$ used to classify students into groups measures some true underlying ability parameter $\theta_i$ with error. 
$$r_i = \theta_i + u_i$$
Where $u_i \sim \mathcal{N}(0,\sigma_u)$. In our model for test scores, we'll be saying that we care about minimizing the distance between the instructional level and the pupil's actual ability level $|\theta_i - \lambda|$, or $(\theta_i - \lambda)^2$. We just want to minimize loss. We want to minimize loss, and there's a question of whether the instructional level assignments under treatment have lower loss than those prevailing under the control conditions.

Our problem is that we do not observe the available instructional levels $\lambda$ and we do not observe the true abilities $\theta_i$. We just have a noisy measure $r_i$. 

We can observe the relationship between $r_i$ and endline scores $y_i$, which are both noisy measures of some true underlying ability parameter. We can relate these by
$$ \theta_{i} = \theta_{i,-1} + \mu - \gamma (\theta_{i,-1} - \lambda)^2 $$
Where the $-1$ subscript indicates that it's from the lagged period, and $\mu$ indicates the amount of schooling that a student who receives optimal instruction would get from this class. 

$E*[y_i|r_i,Z_j,G_i] = \alpha + \beta r_i + (\theta_{i,-1} - Z_i \lambda)$


The default course level is normalized to zero, and the high level is set externally at $\lambda$. $\lambda$ is unobserved.

We have in mind an education production function whereby the amount of learning produced by a class is declining in the distance between the instructional level and the student's true ability level. So adding some time subscripts, we have a model of human capital production in the high course
$$\theta_{i,t} = \theta_{i,t-1}+h(\theta_{i,t-1} - \lambda)$$
and in the low course
$$\theta_{i,t} = \theta_{i,t-1}+h(\theta_{i,t-1} - 0)$$
where $h'<0$. 


We get to decide who receives 

W




Trying to think of some functional forms for $h$: 
$$h(x) = \alpha-\gamma|x|$$
$$h(x) = \frac{1}{\alpha + \gamma |x|}$$
$$h(x) = \alpha - \gamma ln(|x|)$$
Does it make sense to assume that this is symmetric? Seems like there's a difference between something going over my head and reviewing something I already mostly understand.

The BLUP conditional on a bunch of covariates, including the running variable $r_i$ and the pupil's grade is
$\mu_i=E*[y_i|r_i,g_i,X_i]$

In each experiment, student $i$ type $\phi_i = (r_i,g_i)$ is a tuple composed of the student's running score $r_i$ and grade $g_i$. 

Treatment is an assignment rule that assigns all students with 

\sum \mathbbm{1}\{r_i\in R \}
Define a student's potential peer treatment status as $A_i(t) = $


\end{document}